\chapter{Въведение}

\projectname{} представлява курсов проект, насочен към управление на видео анализи с помощта на компютърно зрение. Практическата идея на системата е потребителят да може да качва локални видеа, да създава задачи за анализ върху тях, да следи изпълнението на тези задачи и да преглежда резултатите от разпознаването на обекти. Темата е избрана така, че да бъде едновременно достатъчно практическа и достатъчно богата архитектурно, за да позволи смислена реализация на пълноценен софтуерен сценарий.

Проектът не е замислен като обикновен CRUD пример. Неговата цел е да демонстрира как в рамките на .NET екосистемата може да се изгради приложение с реален домейн, ясно разделение на отговорностите и преизползване на логика между няколко технологични слоя. В случая това преизползване се проявява в две направления. От една страна едни и същи споделени ViewModel-и обслужват два различни потребителски интерфейса. От друга страна едни и същи абстракции на модела и бизнес логиката работят с две различни системи за управление на бази данни.

Функционално системата обединява няколко последователни сценария. Първо се качва изходно видео, което се съхранява във файловото хранилище на приложението. След това върху него се създава задача за анализ. Тази задача се обработва асинхронно от фонов обработчик, който извиква избрания модул за инференция. При реалния сценарий с YoloDotNet се извличат подбрани кадри, върху тях се изпълнява разпознаване на обекти, резултатите се записват в базата данни и накрая се генерира анотирано изходно видео. Потребителят може да преглежда както историята от анализи, така и детайлите за конкретно изпълнение, включително откритите обекти и получения видео резултат.

Архитектурно проектът е реализиран така, че логиката на представяне да бъде максимално споделена, а специфичните за конкретната технология елементи да останат тънък слой върху общия приложен слой. Именно поради това структурата на проекта не е сведена до механични папки от тип \texttt{Views}, \texttt{ViewModels} и \texttt{Models}, а е организирана в по-реалистични слоеве: \texttt{Core}, \texttt{Application}, \texttt{Infrastructure}, \texttt{Web} и \texttt{Avalonia}. Тази организация позволява по-ясно разделение между домейн модела, приложната логика, инфраструктурните зависимости и визуализацията.

Настоящата документация има за цел да представи както крайния резултат, така и инженерните решения зад него. По-конкретно се разглеждат избраните технологии, условията за използването им, архитектурната организация на проекта, начинът на преизползване на ViewModel-и и на СУБД, преизползваемите механизми за филтриране и визуализация на списъци, както и примерни сценарии на работа. В края са обобщени и основните изводи от разработката, включително трудностите, които се наложи да бъдат преодолени.

Изходният код, README описанието и документацията на проекта са публикувани в GitHub репото \url{https://github.com/Apsie09/CV-Analysis-Hub}.
